\section[Engaging Citizens with Open Government Data]{Engaging Citizens with Open Government Data: The Value of Dashboards Compared to Individual Visualizations}
Für meine Literaturanalyse habe ich den Artikel \textit{„Engaging Citizens with Open Government Data: The Value of Dashboards Compared to Individual Visualizations"} \cite{simonofski2022} von Simonofski et al. ausgewählt, veröffentlicht 2022 in ACM Digital Government: Research and Practice.

\smallskip
\noindent
Die Studie befasst sich mit der Frage, wie Open Government Data (OGD) für Bürger:innen besser nutzbar gemacht werden kann. Auf Basis einer systematischen Literaturrecherche wurden 16 Designprinzipien für OGD-Dashboards entwickelt und anschließend mit 108 Teilnehmenden evaluiert. Nachfolgend stelle ich die zentralen Erkenntnisse vor und zeige deren Relevanz für unser Projekt.

\begin{itemize}

    \item \textbf{Dashboards erhöhen das Verständnis}

    Gut gestaltete Dashboards helfen Bürger:innen, OGD deutlich besser zu verstehen als einzelne, isolierte Visualisierungen.

    $\rightarrow$ Für unsere App bedeutet das: Wir sollten Daten nicht isoliert anzeigen, sondern in einem strukturierten, zusammenhängenden Kontext darstellen.

    \item \textbf{Heterogene Nutzergruppen berücksichtigen}

    Die Studie zeigt, dass OGD-Nutzer:innen sehr unterschiedliche technische Hintergründe haben -- von erfahrenen Expert:innen bis hin zu Gelegenheitsnutzer:innen ohne Vorkenntnisse.

    $\rightarrow$ Das spiegelt sich direkt in unseren Personas wider: Von Benjamin als technikaffinern Studierenden bis zu Luca als Gelegenheitsnutzer muss unsere App für alle verständlich sein.

    \item \textbf{Klare Beschriftungen und Erklärungen}

    Ein häufiges Problem bei OGD-Plattformen ist fehlendes Kontextwissen: Nutzer:innen verstehen nicht, was die Daten bedeuten oder woher sie stammen.

    $\rightarrow$ Unsere App muss jeden Datensatz mit verständlichen Beschriftungen und Quellenangaben versehen, damit auch nicht-technische Nutzer:innen die Daten einordnen können.

    \item \textbf{Geeignete Visualisierungsformen wählen}

    Nicht jede Visualisierungsform eignet sich für jeden Datensatz. Die Studie betont, dass die Wahl der richtigen Darstellungsform entscheidend für das Verständnis ist.

    $\rightarrow$ Da wir eine breite Bandbreite an unterschiedlichen Wiener Open Datensätzen visualisieren, müssen wir für jeden Datensatz die passende Visualisierungsform wählen.

    \item \textbf{Datenqualität transparent machen}

    Bürger:innen müssen darauf vertrauen können, dass die angezeigten Daten korrekt und aktuell sind. Fehlende Transparenz über die Datenqualität ist ein zentrales Hindernis für die Nutzung von OGD.

    $\rightarrow$ Besonders relevant im Hinblick auf unsere negative Persona: Unsere App muss sicherstellen, dass Daten aus vertrauenswürdigen Quellen stammen und nicht für manipulative Zwecke missbraucht werden können.

\end{itemize}

\smallskip
\noindent
\textbf{Zusätzliche Einordnung der Studie:}

\smallskip
\noindent
Die Studie basiert auf einer systematischen Literaturrecherche, aus der insgesamt 16 Designprinzipien für OGD-Dashboards abgeleitet wurden. Diese Prinzipien umfassen unter anderem die Berücksichtigung unterschiedlicher Nutzergruppen, die Auswahl geeigneter Visualisierungsformen sowie die Bereitstellung von Kontext- und Qualitätsinformationen.

\smallskip
\noindent
Darüber hinaus wurde eine experimentelle Evaluation mit 108 Teilnehmenden durchgeführt, in der ein Dashboard mit einzelnen Visualisierungen verglichen wurde. Die Ergebnisse zeigen, dass gut gestaltete Dashboards hinsichtlich Verständlichkeit, Usability und wahrgenommener Datenqualität besser abschneiden als isolierte Einzelvisualisierungen.

\smallskip
\noindent
Für unser Projekt sind diese Ergebnisse besonders relevant, da sie verdeutlichen, dass nicht nur die Visualisierung selbst entscheidend ist, sondern vor allem auch die strukturierte Darstellung, ergänzende Kontextinformationen und eine konsequente Nutzerorientierung. Gerade im Bereich offener Verwaltungsdaten ist dies wichtig, um Informationen für unterschiedliche Nutzergruppen verständlich und zugänglich aufzubereiten.
